% !TeX root = ../Tex/main.tex

In this paper I wish to clarify a contradiction between two similar situations that game theory formalises in different ways.
I will first describe the two situations so that the contradiction appears clear to the reader. Safina’s reflections on animal consciousness highlight game theory’s contrasting formalisation of similar situations.

\vspace{1em}
Carl Safina is a biologist and wrote a world-famous essay (\cite{Safina2015}) about animals' minds. The curious thing about this book is that he tries to bring to readers' attention a problem quite close to game theory. He argues that until now, biologists have refused to deepen the study of animals' minds and “souls” and that is high time they did that. Through some examples he provides empirical evidence of behaviours that may make us think of some animals as human or quasi-human. In these examples animals seem to be driven by rationality and emotions, as much as humans are.
In order to convince his colleagues, he infers what may happen in animals mind, what processes go on, what stimuli animals' minds perceive and what drives their actions.

\vspace{1em}
Case 1:
A staff member of a wildlife park in Africa reports an “attack” by an elephant on a tourist. The scene described clearly highlights the signals that the elephant sends to the tourist and clearly identifies the reasons why it is logical to believe that the elephants intentionally avoided hurting them. We can formalise this game though a signalling game in which the elephant sends a signal to threaten the tourist but doesn’t have any intention of hurting it.
In particular the staff member reports signs on the ground that undoubtedly identify the intention to halt the charge against the tourist moments before impact. This first case can be understood as a signalling game, in which the elephant’s display functions as a credible threat rather than an intention to harm.

\vspace{1em}
Another situation, recorded in the same environment, shows how similar conditions can produce an opposite pattern of behaviour.

\vspace{1em}
Case 2:
In the same wildlife park in Africa, a rescue group reached a shepherd who was badly hurt by an elephant. When the staff reached him, they reported that the shepherd stopped them from harming the animal, arguing that the elephant had kept him safe during the night. The reason is that from the behaviour of the animal it appeared evident what was going on in his mind. The elephant soon realised it had misinterpreted the signals from the shepherd, and after hurting him, the animal took care of him until the rescue group reached the place of the accident.
Here, the same signalling structure breaks down: the equilibrium is reached and the elephant attacks. However,
the event suggests a richer model of rationality, where the animal is able to recognise an "error".
\vspace{1em}

Taken together, these cases expose a conceptual gap between real-world behaviour and the assumptions of rational-agent theory. \textit{\color{red}OR } Both scenarios can be framed within game theory, yet each reveals elements of cognition and emotion that standard models overlook.

\vspace{1em}
I believe that the game theory analysis accurately represents most of the aspects of both cases. However, we are still missing something. To describe this situation more accurately, we should introduce a new feature to the way we represent game theory agents. Both animals and humans may be depicted as rational agents, since their behaviour appears to be consistent with the environment and with a broader goal, if not with an aim in mind. Both appear to be able to communicate with each other, indicating that both species are adapting to signals from others and that they can act as informed agents. And both seem to be aware of each other's strengths, at least partially (\cite{schelling_strategy_1960, smith_logic_1973}). Again, in both cases, the animal seems "pacific" and willing to avoid or prevent conflict, at least at a certain point in time.
Why do we see two different behaviours then? I believe these parallels also reveal what is absent from our theoretical framework.

\vspace{1em}
A possible explanation is that the two human beings differed from each other in some feature that was relevant to the context. On the other side, the same explanation may well be applied to the animals, which may be considered to have their own "personalities", as Safina well explains in his book.

\vspace{1em}
This kind of reasoning naturally invites comparison with frameworks that model how hidden factors influence observed outcomes.
Causal inference (\cite{pearl_causality_2009}) suggests that in order to answer this question we may also think of a "variable" that we are mistakenly excluding from the framework. My guess is that this missing variable exists. I will provide my explanation through the means of game theory.

\vspace{1em}
In both cases analogies between men and animals appear to be crystal clear. Most of the aspects they differ are never represented in game theory formalisations. Additionally, game theory is used to study nations “behaviours” often, and this causes little scepticism to scholars. Under the descriptive limits of these games and their assumptions, I believe no harm is caused when I use the terms human, elephant and agent as synonyms.
I will therefore continue with these examples in the belief that they will shed light on the way we theorise a broader set of agents and interactions. These examples serve both to reinforce the empirical evidence presented by Safina and to challenge the theoretical assumptions underlying game-theoretic models.

\vspace{1em}
In doing so, the paper aims to bridge empirical observation and theoretical formalisation in the study of rational agency.